% Transcription — fonds Alexandre Grothendieck, shelfmark 36, batch 2 (pages 21-40).
% Established from the facsimile archives/batches/36/batch-02.pdf (a copy of
% 36.pdf fetched through the project's relay, grothendieck-relay.onrender.com;
% PDF page k = archivists' page 20+k, no cover sheet), per the skill
% transcribe-grothendieck. The folder is 84 pages in five batches, done in
% parallel.
%
% Pass: Opus 5.5 (claude-opus-5-5), 2026-09-26 — first pass, unchecked against the pages by a human.
%
% Pages skipped: 37 (upside-down typed verso of an unrelated Bourbaki draft on
% algebras, nothing in his hand). Page 34 is a blank sheet carrying only his
% title « Picard et dualité » (underlined): a cover, no \page marker, his words
% become the section heading with a \note.
% Pages noted only (published typescript, not re-set): 29, 31, 33.
% Pages transcribed: 21-28, 30, 32, 35, 36, 38-40.
%
% Typescripts in this batch — two, with opposite decisions.
% (1) Pages 21-26: the continuation of the typescript « Le théorème de
% Griffiths par voie algébrique » begun in batch 1 (pp. 9-20), typed pp. 13-18
% (hand-numbered top right). transcripts/36/batch-01.fr.tex exists and decided
% it is an UNPUBLISHED typescript of his, transcribed in full; aligned with
% that decision. Further evidence on these pages: first person (« Nous avons
% utilisé le langage des motifs... », « J'ai donné notes à Berthelot » in ink),
% « Remords à 7.3 », « pas même Lefschetz vache », « par l'âne qui trotte »;
% the numbering (Remarques 7.4, 7.7, Lemmes 7.5, 7.6, 7.8) continues batch 1's
% section 7. It matches none of his published SGA 7 exposés (I, II, VI-IX:
% checked against the tables of contents; Griffiths' theorem is Katz's
% Exposé XX, whose sections do not match this text). The ink additions are
% given with \add{} and notes; machine x-outs not recorded.
% (2) Pages 29, 31, 33: three leaves of the typescript of SGA 7 Exposé VIII
% (« Compléments sur les biextensions », §2 « Accouplements définis par une
% biextension »: 2.1-2.2.3 on p. 33, 2.2.4-2.3.3 on p. 31, end of 2.3 on p. 29),
% scanned upside down as the backs of pp. 28, 30, 32. Evidence: section title
% and numbering (2.1.2), (2.2.2), (2.2.6), (2.3.1)-(2.3.5) coincide with the
% published exposé. PUBLISHED text: not re-set; one French \note per page
% identifying it, with the ink interventions (presumably his) in full. Their
% top-right numerals are not read as a pagination (they are the reverse
% show-through of the archivists' numbers).
%
% His pagination: typed pp. 13-18 on pp. 21-26 (the Griffiths typescript ends
% on p. 26 with an inked heading « 8. Relations heuristiques entre cohomologie
% cristalline ... » and no text). No pagination of his on the manuscript notes
% 27-40. Numbered runs: Remarques 7.4 a)-b), 7.7 a)-c), Lemmes 7.5, 7.6, 7.8;
% (*), (**), (***) on pp. 22-23; (*), (**) on p. 38.
%
% What the batch is about. (a) End of the Griffiths typescript: Remarques 7.4
% (7.3 without n even, level of E^{n-1}(Y_t) is n-1; stated without motives);
% proof of 7.1 via Lefschetz-Verdier (Deligne) congruence card Y_s(k') = c +
% (-1)^{n-1} Tr f_s mod p, Lemmes 7.5-7.6 (forms without non-trivial zero,
% normic forms, estimate of d'_0), Remarques 7.7, Lemme 7.8 (specialisation of
% the semisimple part of a p-linear endomorphism). (b) Manuscript notes on
% biextensions of complexes [M -> G] (G extension of B by a torus D(M')), the
% l-adic pairing on T_l(A) x T_l(A'), orthogonality of D(M'_n) and nG'
% (pp. 27-32). (c) « Picard et dualité »: a table of H^i(X, z_l), H^i(X, z_p),
% H^i(X, i_p) for a curve (p. 35); Ext_X(f*B, G_m) vs Ext_S(B, f^!G_m) and a
% list of topics (p. 36); H^1(X, A) = Ext^1(A', Pic_X) for an abelian variety
% A, compared with H^1(X, G) = Hom(D(G), Pic_X) for G finite (pp. 38-39);
% relative H^1(X/S; Hom(gamma, G)) = Hom(gamma, H^1(X/S; G)) (p. 40).
%
% Notation: \ell for the typewriter chi; \phi, \underline{\phi}, \psi for the
% typed ø, underlined ø and slashed ø on p. 24; \mathfrak{z}_\ell for his
% curly z (p. 35); \gamma and \underline{\mathcal{H}} on p. 40.
%
% Apparatus: 52 \ill{}, 80 \uncertain{}. Hardest
% pages: 28, 30 and 32 (fast notes, much of the connective prose lost), 36
% (the list of topics), and the last lines of 39. Readings to check first: the
% « Exemple » of p. 28; the opening of p. 30; « Berthelot » on p. 26; the
% marginal of p. 39.

\documentclass[11pt,a4paper]{article}
% Le préambule commun à toutes les transcriptions du fonds Grothendieck.
%
% Il tient en une idée : **l'incertitude se compose, elle ne se résout pas.**
% Une transcription qui lisse un mot douteux a détruit la seule chose pour
% laquelle on la faisait — un lecteur doit pouvoir savoir, à chaque endroit,
% ce qui était sur le feuillet et ce qui a été ajouté. Les six macros qui
% suivent sont donc l'appareil critique, et non de la décoration.
%
% Les mêmes commandes sont reconnues par `scripts/render.mjs`, qui produit la
% vue de lecture du site. Une macro ajoutée ici doit l'être aussi là-bas,
% faute de quoi le rendu échouera bruyamment — ce qui est voulu : un rendu
% silencieusement faux est pire qu'un rendu absent.

\ProvidesPackage{grothendieck}[2026/08/08 Transcriptions du fonds Grothendieck]

\usepackage{amsmath,amssymb,amsthm}
\usepackage{mathrsfs}
\usepackage{stmaryrd}
% Les diagrammes commutatifs sont partout dans ce fonds : c'est la langue
% même de la géométrie algébrique de Grothendieck, et une transcription qui
% les rendrait en prose ne transcrirait rien.
\usepackage{tikz-cd}
% Français par défaut — c'est la langue des feuillets — mais l'anglais est
% chargé aussi : la traduction et le résumé sont des documents anglais, et la
% typographie française (espaces avant les deux-points) y serait une faute.
% Ces documents basculent par \selectlanguage{english} en tête de corps.
\usepackage[english,french]{babel}
\usepackage{fontspec}
\usepackage{geometry}
\geometry{a4paper,margin=25mm,marginparwidth=28mm}
\usepackage{marginnote}
\usepackage{xcolor}
% ulem, not soul. `soul`'s \st and \ul are built on a character-by-character
% scan that XeLaTeX's font machinery breaks: the document fails with an
% undefined control sequence rather than degrading. ulem does the same two jobs
% and is Unicode-engine safe. `normalem` stops it hijacking \emph, which the
% transcriptions use for Grothendieck's own emphasis.
\usepackage[normalem]{ulem}
\usepackage{hyperref}
\hypersetup{colorlinks=true,linkcolor=black,urlcolor=black,pdfborder={0 0 0}}

% --- Métadonnées du lot -----------------------------------------------------
% Reprises telles quelles par le script de rendu, qui les affiche en tête de
% la vue de lecture. Elles fixent ce que le fichier prétend couvrir : sans
% elles, une transcription flotte sans point d'attache dans un fonds de seize
% mille feuillets.

\newcommand{\folder}[1]{\gdef\grothFolder{#1}}
\newcommand{\batch}[1]{\gdef\grothBatch{#1}}
\newcommand{\leaves}[2]{\gdef\grothFirst{#1}\gdef\grothLast{#2}}
\newcommand{\foldertitle}[1]{\gdef\grothFoldertitle{#1}}
\gdef\grothFolder{?}\gdef\grothBatch{?}\gdef\grothFirst{?}\gdef\grothLast{?}\gdef\grothFoldertitle{}

% --- Le résumé --------------------------------------------------------------
%
% Il ouvre la lecture modernisée, sous le titre « Résumé », et il est composé
% en petit. Ce n'est pas un ornement : c'est le seul passage du document qui
% ne soit pas des mathématiques, et un lecteur doit pouvoir le reconnaître
% avant de l'avoir lu — puis le sauter s'il connaît déjà le sujet, ou s'y
% arrêter s'il ne le connaît pas. Le filet qui le suit marque la reprise du
% corps.

\newenvironment{resume}
  {\small\setlength{\parskip}{0.45em}}
  {\par\medskip\hrule\medskip}

% --- Mots-clés ---------------------------------------------------------------
%
% En anglais, en clôture du résumé de la lecture modernisée : le vocabulaire
% moderne sous lequel chercher ce que les feuillets construisent. Le manifeste
% du site (`npm run manifest`) les extrait pour étiqueter la cote entière —
% cette ligne est la source unique des tags, il n'y a pas de fichier à côté.
\newcommand{\keywords}[1]{\par\smallskip\noindent
  {\footnotesize\textsc{Keywords} --- \textit{#1}}\par}

% --- L'appareil critique ----------------------------------------------------

% Le numéro de feuillet, en marge. C'est l'ancre qui permet de régler un
% désaccord sur une formule en regardant *un* feuillet plutôt qu'une chemise
% de six cents. Le site s'en sert aussi pour tourner les pages du fac-similé
% quand on fait défiler la transcription.
\newcommand{\leaf}[1]{%
  \marginnote{\footnotesize\color{gray!70}\textsf{#1}}%
  \label{leaf:#1}\ignorespaces}

% Illisible : on ne devine pas. Un mot inventé qui se lit comme les autres est
% le pire résultat possible.
\newcommand{\ill}{\textcolor{red!60!black}{[\,\ldots\,]}}

% Lecture douteuse : le mot est donné, et signalé comme incertain.
\newcommand{\uncertain}[1]{\uline{#1}}

% Ajout de l'éditeur — une lettre manquante, un mot sous-entendu.
\newcommand{\add}[1]{\textcolor{blue!50!black}{[#1]}}

% Biffé par Grothendieck lui-même. Ce qu'il a rayé fait partie du document :
% une rature dit souvent où la pensée a changé de direction.
\newcommand{\struck}[1]{\sout{#1}}

% Note du transcripteur, sur l'état matériel du feuillet.
\newcommand{\note}[1]{\par\smallskip{\small\color{gray!60!black}\textit{[#1]}}\par\smallskip}

% Note marginale de Grothendieck, distincte du corps du texte.
\newcommand{\marginal}[1]{\par\smallskip\noindent\hspace{1em}%
  {\small\color{gray!60!black}#1}\par\smallskip}

% --- Titre ------------------------------------------------------------------

\AtBeginDocument{%
  \begin{center}
    {\large\bfseries Fonds Alexandre Grothendieck --- cote n\textsuperscript{o}~\grothFolder}\\[2pt]
    {\itshape\grothFoldertitle}\\[4pt]
    {\small feuillets \grothFirst--\grothLast\ (lot \grothBatch)}\\[2pt]
    {\footnotesize\color{gray!60!black}Transcription établie d'après le fac-similé numérisé
     par l'Université de Montpellier.\\
     Les lectures incertaines sont soulignées, les illisibles notées [\,\ldots\,],
     les ajouts entre crochets.}
  \end{center}
  \vspace{1em}}

\endinput


\folder{36}
\batch{2}
\pages{21}{40}
\dating{[vers 1967-1973]}
\watermark{Édition de démonstration}
\foldertitle{SGA 7 (ma part) : notes manuscrites (s.d.), tapuscrit annoté (s.d.)}

\begin{document}

\section*{Le théorème de Griffiths par voie algébrique (suite)}

\note{suite du tapuscrit commencé aux pages 9 à 20 ; le titre de section est
repris de la page 9}

\page{21}

\note{p.~13 du tapuscrit (numéro à la main en haut à droite) : la suite
immédiate de la p.~12 (page 20), sur laquelle s'arrête le lot précédent}

\emph{Remarques} 7.4.\quad a) Remords à 7.3. Inutile de supposer $n$ pair et
$n \geqslant 4$, en énonçant la conclusion sous la forme : le niveau de
$E^{n-1}(Y_t)$ est égal à $n-1$ ($n \geqslant 1$). On n'a pas besoin en effet
d'un critère de Ogg-Néron-Chafarévitch ici, car il suffit de savoir que
\emph{si} $E^{n-1}(Y_t)$ est de niveau $< n-1$, i.e.\ s'il existe sur
$\overline{k'}$ ($k' = k(t)$) un cycle sur un produit $Z \times
Y_{\overline{k'}}$ définissant un \add{épi}morphisme $H^{n-3}(Z,
\underline{\mathbf{Q}}_\ell)(-1) \to H^{n-1}(Y_{\bar{k}'},
\underline{\mathbf{Q}}_\ell) \to E(Y_{\bar{k}'})$, alors pour tout $s \in D$
assez voisin de $t$, la même conclusion est vraie pour $E(Y_s)$, comme on
constate aussitôt, de sorte que, si $s$ est à corps résiduel fini, les valeurs
propres du frobenius opérant sur $E^{n-1}(Y_{\bar{s}})$ sont divisibles par card
$k(t)$ et a fortiori par $p$. Or cela est impossible si $D$ est choisie assez
générale, de façon qu'on puisse appliquer Bertini. D'autre part, même si $n$
est impair i.e.\ $n-1$ pair, de sorte qu'on ne peut plus affirmer que la
représentation de monodromie sur $E^{n-1}(Y_{\bar t})$ reste simple quand on la
restreint à un sous-groupe ouvert quelconque de $G =
\operatorname{Gal}(\overline{k(t)}/k(t))$, il sera vrai néanmoins qu'aucun
sous-motif non nul de $E^{n-1}(Y_{\bar t})$ n'est de niveau $< n-1$. En effet,
le plus grand sous-motif \add{de niveau $< n-1$} de $E^{n-1}(Y_{\bar t})$
définit un sous-espace de $E^{n-1}(Y_{\bar t}, \underline{\mathbf{Q}}_\ell)$,
qui est évidemment stable sous le groupe de Galois.\note{« épi » est tapé dans
l'interligne au-dessus de « morphisme », sur un préfixe barré à la machine ;
de même « de niveau $n-1$ » (sic, pour $< n-1$ ?) est tapé au-dessus de « de
$E^{n-1}$ » et appelé par un trait ; les passages effacés à la machine (« k(t) »
avant $\bar{k}'$, une ligne entière avant « alors pour tout », une ligne avant
« représentation ») ne sont pas repris. Le $\ell$ de la machine est, comme dans
le lot précédent, un caractère en forme de $\chi$}

b)\quad Nous avons utilisé le langage des motifs par pure commodité pour
l'énoncé de 7.3. Comme ce langage n'est pas justifié à l'heure actuelle, il est
préférable, pour le moment, de l'énoncer en termes de la filtration par le
niveau sur la cohomologie $\ell$-adique, en définissant celle-ci en conformité
avec a) (La définition plus naturelle en termes de la filtration de l'espace
$Y_{\bar t}$ par les supports devrait être équivalente, mais ceci utiliserait
la résolution des singularités). Ceci dit, l'énoncé 7.3 (sous la forme plus
générale a) ci-dessus) est démontré sans aucune conjecture (pas même Lefschetz
vache pour $X$ \ldots).

\page{22}

\note{p.~14 du tapuscrit. Plusieurs lignes de ce feuillet et des suivants sont
doublées sur la reproduction, un peu décalées ; on ne les lit qu'une fois}

\emph{Démonstration de} 7.1.\quad Nous prendrons \struck{pour} $d_0$
\struck{le plus petit entier} tel que l'on ait \add{(entre autres)}
\[
(*) \qquad H^i(X, \underline{\mathcal{O}}_X(-d)) = 0
\quad \text{pour tout } d \geqslant d_0,\ i < n.
\]
\note{« (entre autres) » est à l'encre ; « le plus petit entier » est barré à la
machine, « pour » à l'encre ; entre « (*) » et la formule, un début de formule biffé à l'encre,
\struck{$d_0$ \ldots{} $n+1$} ; les signes « (*) » et « (**) » en marge sont à la main}
Désignons par $U$ l'ouvert de $S_d$ formé des $s \in S_d$ tels que $\dim Y_s =
n-1$, qui est aussi le plus grand ouvert tel que $\underline{Y} \to S_d$ soit
plat au dessus de $U$. Comme on a signalé, le choix de $d_0$ implique que pour
tout $s \in U(k')$, les homomorphismes
\[
(**) \qquad H^i(X_{k'}, \underline{\mathcal{O}}_{X_{k'}}) \longrightarrow
H^i(Y_s, \underline{\mathcal{O}}_{Y_s})
\]
sont des isomorphismes pour $i \leqslant n-2$, un monomorphisme pour $i = n-1$.
Ceci implique, par EGA III 7, que les $R^i f_{*}(\underline{\mathcal{O}}_Y)$ sont
plats sur $U$ et y commutent à l'extension de la base, en particulier au passage
aux fibres\add{.} On fera attention que même si $X$ est lisse et géom.\ connexe,
nous serons obligés de façon essentielle de travailler avec des points $s \in U$
pour lesquels $Y_s$ n'est pas irréductible et a fortiori non lisse. Ce ne sera
qu'une fois que 7.1 est démontré qu'on pourra en tirer des conclusions
diophantiennes sur des points fermés dans un ouvert arbitrairement petit de $S_d$
comme dans 7.2, puis les appliquer comme dans 7.3.\note{le « $\leqslant$ » de
« $i \leqslant n-2$ » est surchargé à la main ; « diophantiennes » est tapé dans
l'interligne}

Bien entendu, il suffit de prouver 7.1 après extension arbitraire du corps de
base $k$, ce qui nous permettrait en particulier de supposer $k$ algébriquement
clos ; et la conclusion se renforce \add{dans le cas général} automatiquement, en
remplaçant le groupe de Galois $G$ envisagé par le sous-groupe image du groupe
de Galois provenant de la situation géométrique où $k$ est remplacé par
$\bar{k}$. Mais pour prouver 7.1, nous procèderons \add{au contraire} par voie
arithmétique, en notant qu'on peut se réduire au cas où $k$ est fini. Cette
réduction se fait de la façon habituelle, et n'offre pas de difficultés. On peut
alors considérer les frobenius « géométriques » opérant sur les $H^i(X,
\underline{\mathcal{O}}_X)$ ; si $k$ a $q = p^m$ éléments, \struck{c'est}
\add{ce sont} \uncertain{les} puissance\add{s} $m$-ème\add{s} \uncertain{des}
frobenius\note{« dans le cas général » est tapé dans l'interligne et appelé par
un trait ; « au contraire » est à l'encre ; à la dernière ligne, l'encre
corrige « c'est la puissance $m$-ème du » en « ce sont les puissances $m$-èmes
des » (« sont » au-dessus de « est » barré, surcharges sur « la » et « du »)}

\page{23}

\note{p.~15 du tapuscrit}

$p$-linéaire\add{s} opérant sur les $H^i(X, \underline{\mathcal{O}}_X)$. Quitte
à remplacer $k$ par une extension finie, on peut supposer que toutes les valeurs
propres desdits frobenius sont égales à $0$ ou $1$. Soit $c$ la somme alternée
des rangs semi-simples des $H^i(X, \underline{\mathcal{O}}_X)$ pour $0 \leqslant
i \leqslant n-1$\add{.} Alors pour toute extension finie $k'$ de $k$, la somme
alternée des traces de $\mathrm{frob}_{k'}$ opérant sur les $H^i(X_{k'},
\underline{\mathcal{O}}_{X_{k'}})$ est \struck{alors} égale à $c$. En vertu de
Lefschetz-Verdier Deligne, et des isomorphismes $(**)$, on conclut que pour tout
point $s \in U(k')$, on a
\[
(***) \qquad \operatorname{card} Y_s(k') \equiv c + (-1)^{n-1} \operatorname{Tr} f_s ,
\]
où $f_s \in G_0$ est le frobenius relatif à $k'$ opérant sur $E_0(Y_s)$. Ceci
posé, 7.1 est une conséquence des deux lemmes suivants ;\note{le « s » de
« $p$-linéaires » est à la main ; « des valeurs » et une ligne entière avant
« des rangs semi-simples », puis « l'image de l'élément de Frobenius de »
avant « le frobenius relatif », sont effacés à la machine ; « alors » est barré
à l'encre ; le « (***) » de la marge et l'indice $s$ de $Y_s$ sont à la main}

\emph{Lemme} 7.5.\quad Pour tout $s \in U$, il existe un élément $f \in G_0$
\add{et $V'_0 \subset V_0$ \uncertain{stable par} $f$} tel que le polynôme
caractéristique de $f$ opérant sur $V'_0$ soit égal à celui de $f_s$, et en
particulier tel que $f$ et $f_s$ aient même trace.\note{l'ajout « et $V'_0
\subset V_0$ stable par $f$ » est à l'encre au-dessus de la ligne et appelé par
un trait après $G_0$ ; le $'$ de $V'_0$ est ajouté à la main et $V'_0$ entouré ;
de même un $'$ sur le $f$ de la dernière ligne, en partie effacé. Un trait
vertical à l'encre borde l'énoncé dans la marge gauche}

\emph{Lemme} 7.6.\quad Pour tout entier $N \geqslant 0$, il existe un entier
$d'_0$ \add{$> 0$} tel que pour tout entier $d \geqslant d'_0$, il existe une
extension finie $k'$ de $k$ et un point $s \in U_d(k')$ tel que
$\operatorname{card}(Y_s) = N$. (Marche dès que $X$ est un sous-schéma fermé de
$\mathbf{P}^r_k$ de dimension $n \geqslant 1$.)\note{« $> 0$ » est à l'encre}

\emph{Démonstration de} 7.6.\quad \add{\uncertain{On peut supposer $X$
réduit.}} Soit $r$ la \add{somme des} degré\add{s} des composantes irréductibles
réduites de $X$ qui sont de dimension $n$, et soit $d_1$ un entier tel que
\[
N \geqslant r d_1^{\,n} .
\]
Quitte à passer à une extension finie de $k$, on peut supposer les composantes
irréductibles $X_i$ de $X$ géométriquement irréductibles, que tout $X_i$ contient
un $x_i \in X_i(k)$, et qu'il existe $n$ formes $\phi_0, \ldots, \phi_{n-1}$ de
degré $d_1$ sur $\mathbf{P}^r_k$\add{,}\note{« On peut supposer $X$ réduit. »
est à l'encre, au-dessus du début de la phrase, qu'une parenthèse ouvre ;
« somme des » est tapé dans l'interligne, et « la degré » corrigé en « la somme
des degrés » ; trois lignes effacées à la machine précèdent « les composantes
irréductibles » (un premier essai sur une composante irréductible $X_1$ de $X$
et un point rationnel)}

\page{24}

\note{p.~16 du tapuscrit. Les formes sont notées à la machine par un « ø » :
$\phi_0, \ldots, \phi_n$ ; la forme à $n+1$ variables est un « ø » souligné,
rendu $\underline{\phi}$, et la forme de degré $d_3$ un « ø » barré d'une
oblique, rendu $\psi$}

définissant des hypersurfaces $H_0, \ldots, H_{n-1}$ telles que $X \cap H_0 \cap
\ldots \cap H_{n-1} = Z$ soit un schéma fini \add{étale} sur $k$, de degré $r
d_1^{\,n}$, dont tous les points sont rationnels sur $k$ \add{et distinct des
$x_i$}. Quitte à augmenter $k$ et $d_1$, on peut supposer qu'il existe une forme
$\phi_n$ de même degré $d_1$, définissant $H_n$, telle que $H_n \cap Z$ soit de
cardinal $N$. Pour tout entier $d_2 \geqslant n+1$, on sait qu'il existe une
forme $\underline{\phi}$ à $n+1$ variables définie sur $k$, de degré $d_2$,
n'ayant pas de zéro non trivial sur $k$. Alors
\[
\phi = \underline{\phi}(\phi_0, \ldots, \phi_n)
\]
est une forme de degré $d_2 d_1$ sur $\mathbf{P}^r_k$, telle que $(H_\phi \cap
X)(k) = (H_n \cap Z)(k)$ soit de cardinal $N$. De plus, comme cet ensemble ne
contient pas les $x_i$, $\phi$ définit bien un point de $U_{d_1 d_2}$. Lorsque
$N$ est $0$ ou $1$ \add{(plus généralement $N \leqslant r$)}, on a terminé, car
on peut prendre $d_1 = 1$, et 7.6 marche avec $d'_0 = n+1$. Si on a $N > r$, on
note que pour tout $d_3 \geqslant n+1$, il existe une forme $\psi$ de degré
$d_3$ \add{rat/$k$} telle que $(H_\psi \cap X)(k) = \emptyset$ d'après ce qu'on
vient de voir, et alors la forme $\phi\psi$ est de degré \add{$d =$} $d_1 d_2 +
d_3$ et $(H_{\phi\psi} \cap X)(k) = Z(k)$ a comme cardinal $N$, et $\phi\psi$
définit encore un élément de $U_d$. On peut donc prendre alors
\[
d'_0 = d_1 d_2 + (n+1) = (d_1 + 1)(n+1)
\]
(en prenant, comme il est loisible, $d_2 = n+1$.)\note{« étale », « et distinct
des $x_i$ », « (plus généralement $N \leqslant r$) », « rat/$k$ » et « $d =$ »
sont des ajouts interlinéaires à l'encre appelés par des traits. Tel quel, le
texte écrit $(H_\phi \cap X)(k)$, $(H_{\phi\psi} \cap X)(k) = Z(k)$ et $\phi\psi$
« de degré $d_1 d_2 + d_3$ » : on ne corrige pas}

\emph{Remarques} 7.7.\quad a) On trouve ainsi une estimation assez précise du
$d'_0$ de 7.6 en fonction de $N$ et de $r$, donc du $d_0$ de 7.1 en fonction de
$N = p-1$ et de $r$, et du plus petit entier rendant valables les égalités $(*)$.
Lorsque dans la conclusion de 7.1 on se contente d'affirmer que les éléments de
$G_0$ donnent lieu à au moins deux traces distinctes (ce qui suffit à conclure
que $G_0$ n'est pas un $p$-groupe et a fortiori n'est pas le groupe unité, donc
$\dim V \geqslant 1$ i.e.\ frobenius sur $V$ n'est pas nilpotent) on peut prendre
simplement $d'_0 = n+1$, et pour $d_0$ le plus petit entier $\geqslant n+1$ tel
que l'on ait $(*)$. Ce choix de $d_0$ est suffisant aussi, par conséquent, pour
avoir la validité des\note{le « a) » est repassé à la main}

\page{25}

\note{p.~17 du tapuscrit}

conclusions de 7.2\add{,} sous la forme affaiblie qu'on peut satisfaire à la
congruence envisagée pour au moins deux valeurs de $a$ distinctes mod.~$p$, et
par suite pour avoir la conclusion de 7.3 (qui n'utilise essentiellement que 7.2
sous cette forme affaiblie).

b)\quad Lorsqu'on veut 7.6 avec un $N$ donné, on pourra sans doute trouver
souvent une meilleure majoration de $d'_0$ que $(d_1 + 1)(n+1)$, en décomposant
$N$ en $N' + N''$, en choisissant $d'_1$ et $d''_1$ de façon que $N' \geqslant r
d_1'^{\,n}$, $N'' \geqslant r d_2'^{\,n}$, et construisant une forme produit de
degré \add{$d =$} $d'_1 d'_2 + d''_1 d''_2$ ($d'_2, d''_2$ des entiers
arbitraires $\geqslant n+1$) dont l'ensemble des zéros communs avec $X$
rationnels sur $k$ soit la somme disjointe de deux ensembles de cardinaux
respectifs $N'$ et $N''$. Quitte à choisir $d'_1$ et $d''_1$ premiers entre eux,
on trouve \add{que} tout entier assez grand $d$ \add{se met} sous la forme
envisagée.\note{un crochet à l'encre sépare « b) » de ce qui précède, et le
« b) » est à la main ; « $d =$ », « que » et « se met » sont à l'encre dans
l'interligne, ce dernier au-dessus de « sous » ; le tapuscrit écrit « $N''
\geqslant r d_2'^{\,n}$ » (sic, pour $d''_1$ ?)}

c)\quad On fera attention que les formes auxiliaires sans zéro non trivial
utilisées dans la démonstration de 7.6 sont construites avec des formes
normiques, et sont toujours décomposables géométriquement en facteurs linéaires.
Ceci précise la remarque faite au début de la démonstration de 7.1 concernant la
nécessité de travailler avec le gros ouvert $U_d$ tout entier.

Indiquons pour terminer la raison générale pour 7.5 :

\emph{Lemme} 7.8.\quad Soient $S$ un schéma \add{normal} de car.\ $p > 0$,
$\underline{V}$ un Module localement libre de type fini sur $S$ muni d'un
\add{endo}morphisme $p$-linéaire $f$, $t$ un point de $S$ et $s$ une
spécialisation de $t$. Considérons les représentations des groupes $G_t =
\operatorname{Gal}(\overline{k(t)}/k(t))$ et $G_s =
\operatorname{Gal}(\overline{k(s)}/k(s))$ dans des vectoriels de dimension finie
$V_1$ resp.\ $V_0$ sur le corps premier $\mathbf{F}_p$, définis par les parties
semi-simples des fibres réduites de $\underline{V}$ en $t$ resp.\ en $s$. Alors il
existe un sous-espace $V'_0$ de $V_1$, stable par un sous-groupe ouvert $G'_t$ de
$G_t$, tel \add{que} la représentation de l'image $G'_{t\,0}$ de $G'_t$ dans
$\operatorname{Aut}(V'_0)$ par automorphismes de $V'_0$ soit isomorphe à celle de
l'image $G_{s\,0}$\note{« normal » est écrit dans la marge gauche et appelé
après « schéma » par une oblique ; « endo » et « que » sont à l'encre dans
l'interligne. Le texte tapé porte « point de s » pour « point de $S$ »}

\page{26}

\note{p.~18 du tapuscrit, la dernière : le texte tapé s'arrête au premier tiers
du feuillet, le titre du n\textsuperscript{o}~8 est à l'encre, et le reste du
feuillet est vierge}

de $G_s$ dans $\operatorname{Aut} V_0$ par automorphismes de $V_0$, via un
isomorphisme convenable de $G'_{t\,0}$ avec $G_{s\,0}$.

La démonstration de 7.8 est par l'âne qui trotte, en se ramenant au cas $S$
\add{intègre} hensélien de point fermé $s$, \add{point générique $t$}.\note{« 
intègre » et « point générique $t$ » sont à l'encre, le premier au-dessus de la
ligne, appelé par un trait, le second en bout de ligne, sur le point final
tapé}

8.\quad \marginal{Relations heuristiques entre cohomologie cristalline et la
\uncertain{torsion} cristalline, et la cohomologie cohérente. (J'ai donné notes à
\uncertain{Berthelot})}\note{titre du n\textsuperscript{o}~8 et parenthèse écrits
à la main, à l'encre, sous le texte tapé ; aucun texte ne suit}

\section*{Biextensions et complexes $[M \to G]$}

\note{à partir d'ici, notes manuscrites de sa main, à l'encre, sans pagination
de l'auteur. Le titre de section est de l'éditeur. Les feuillets 28, 30 et 32
sont écrits au recto de feuillets du tapuscrit de l'exposé VIII de SGA~7
(pages 29, 31, 33 ci-dessous) ; le feuillet 27 est une feuille de papier
chamois qui ne porte que le croquis qui suit, en haut à droite}

\page{27}

$M$ \uncertain{Situation}

\begin{tikzcd}[column sep=small, row sep=small]
0 \arrow[r] & T \arrow[r] \arrow[d, no head, "\simeq"'] & G \arrow[r] & B \arrow[r] & 0 \\
 & D(M') & & & \\
0 \arrow[r] & T' \arrow[r] \arrow[d, no head, "\simeq"'] & G' \arrow[r] & B' \arrow[r] & 0 \\
 & D(M) & & &
\end{tikzcd}

\note{une flèche va de $M$ vers $G$ au-dessus de la première ligne, et une autre
de $M'$, écrit sous $D(M)$, vers $G'$ ; les isomorphismes $T \simeq D(M')$,
$T' \simeq D(M)$ sont notés par des $\simeq$ verticaux. Une croix entre les deux
lignes, sous $B$, et des signes \ill{} près de $M'$}

\page{28}

$\Gamma = \mathbb{G}$ \quad $M \times B'$ \quad $\operatorname{Ext}^1(B', DM)$
\quad \struck{$\operatorname{Hom}(B', \operatorname{Ext}^1(M, \mathbb{G}_m)) = 0$}
\note{en haut de la page ; « $\Gamma = \mathbb{G}$ » est à gauche, le reste à
droite ; la dernière formule est biffée}

\[
\begin{array}{cc}
B & B' \\
{\scriptstyle i}\uparrow & \uparrow{\scriptstyle i'} \\
M & M'
\end{array}
\qquad \text{biext.\ } \varepsilon \text{ vers } \mathbb{G}_m
\]

\[
B' \xrightarrow{\ \alpha'\ } \underline{\operatorname{Ext}}^1(B, \mathbb{G}_m) \to
\underline{\operatorname{Ext}}^1(M, \mathbb{G}_m) \ \text{nul}
\qquad \text{p.ex.\ } \operatorname{Hom}(B', \underline{\operatorname{Ext}}^1(M, \mathbb{G}_m)) = 0 \ ?
\]
\[
B \xrightarrow{\ \alpha\ } \underline{\operatorname{Ext}}^1(B', \mathbb{G}_m) \to
\underline{\operatorname{Ext}}^1(M', \mathbb{G}_m) \ \text{nul}
\qquad \text{p.ex.\ } \operatorname{Hom}(B, \underline{\operatorname{Ext}}^1(M', \mathbb{G}_m)) = 0 \ ?
\]
\note{$\alpha'$, $\alpha$ surmontent des $\simeq$ ; le point d'interrogation de
droite est un « ? » suivi d'un guillemet de répétition sous lui}

d'où extensions,
\[
0 \to DM \xrightarrow{\ j'\ } G' \xrightarrow{\ p'\ } B' \to 0, \qquad
0 \to DM' \xrightarrow{\ j\ } G \xrightarrow{\ p\ } B \to 0
\]
\marginal{$\mathcal{E}_{M,G'}$ est canoniqt splitté par $\rho$ ; $\mathcal{E}_{G,M'}$
est can.\ splitté $\rho'$}\note{sous $DM$ est écrit $T'$, sous $DM'$ un $T$ ;
les deux phrases de droite sont séparées du reste par des traits horizontaux}

\emph{Exemple} : si \struck{\ill{}} $\alpha$ \add{\ill{}} un isom., alors il
revient au même de donner un $i \colon M' \to \underline{\operatorname{Ext}}^1(B,
\mathbb{G}_m)$
\[
\operatorname{Hom}(M', \underline{\operatorname{Ext}}^1(B, \mathbb{G}_m))
\rightleftarrows \operatorname{Biext}(B, M'; \mathbb{G}_m) \leftleftarrows
\operatorname{Ext}^1(B, D(M'))
\]
\uncertain{si} $\underline{\operatorname{Ext}}^1(M', \underline{\operatorname{Hom}}(B,
\mathbb{G}_m)) = 0$, \uncertain{vrai} \add{car} $\underline{\operatorname{Hom}}(B,
\mathbb{G}_m) = 0$ ; si $\operatorname{Hom}(B, \underline{\operatorname{Ext}}^1(M',
\mathbb{G}_m)) = 0$, p.ex.\ $\underline{\operatorname{Ext}}^1(M', \mathbb{G}_m) =
0$.\note{les deux conditions sont écrites sous les deux flèches, chacune appelée
par un crochet ; les deux dernières sont entourées. Au-dessus de
$\operatorname{Ext}^1(B, D(M'))$, un mot biffé}

\uncertain{ainsi} \uncertain{Donc} : \uncertain{la donnée} de $i'$ équivaut à celle
de l'ext.\ $G$. Et réciproquement.

Considérons l'image inverse de $\varepsilon$ sur $(M, M')$ \add{$\varepsilon_{i,i'}$}.
Alors des splittings \add{$\rho_i$} de cette biextension correspondent aux homs de
$M$ dans $G$ relevant $i$, ou aussi de $M'$ dans $G'$ relevant $i'$ :
\[
M \to G \to B, \qquad M' \to G' \to B'
\]
\struck{(sur $k$)}
\note{« $\varepsilon_{i,i'}$ » est entouré au-dessus de « $\varepsilon$ » ;
« $\rho_i$ » au-dessus de « splittings »}

\struck{Donnons-nous $f$ de $j'$ et correspondant ainsi \ill{} $\mathcal{E}_{G,M'}$
et un splitt.\ $\sigma$ canonique \ill{} de $\mathcal{E}_{M,G'}$ \ill{} splitting
$\sigma$ \ill{} splitting \ill{} $\mathcal{E}_{M,M'}$}\note{passage biffé de
traits obliques et encadré ; lecture très incertaine}

Les deux splittings de $\mathcal{E}_{M,M'}$, déduits d'une part d'un splitting
$\rho$ de $\mathcal{E}_{M,G'}$ et de $g \colon M' \to G'$, d'autre part d'un
splitting $\rho'$ de $\mathcal{E}_{G,M'}$ et de $f \colon M \to G$,
coïncident.\note{dans « $\mathcal{E}_{G,M'}$ » de la dernière ligne, les indices
sont surchargés}

\page{29}
\note{feuillet dactylographié, au verso de la page 28 (reproduit tête-bêche) :
feuillet non paginé du tapuscrit de l'exposé VIII de SGA~7 (« Compléments sur les
biextensions. Propriétés générales des biextensions des schémas en groupes »,
§2 « Accouplements définis par une biextension »), fin de 2.3 : la
compatibilité symétrique (2.3.4), (2.3.5), et le début de la preuve de (2.3.1)
par $\operatorname{Tor}_1$. Texte publié, non reproduit. Interventions à la
main : au-dessus de « à son tour », « implique » remplace « est » barré, et
« de (2.3.1) » est ajouté après « compatibilité symétrique » ; « (2.2.2) »
corrigé en surcharge ; le numéro de la formule corrigé en « (2.3.4) » et
« (2.3.5) » ; dans le diagramme (2.3.4), les flèches verticales portent à
l'encre « id », « $m$ », « id », les indices $n$, $n'$ sont repris et la
flèche du bas marquée $\phi^n_\xi$ ; « or » au-dessus de « dans » barré, et
« (2.3.3) » au-dessus d'une égalité ; un $\xi$ ajouté en indice aux
$\phi$ ; « en $Q$ » après « fonctorialité » ; une ligne entière de formule
barrée à la machine ; dans le diagramme final, « $m \otimes \mathrm{id}_Q$ »
écrit à la main à côté de la flèche verticale. Main vraisemblablement la
sienne}

\page{30}

\uncertain{Lorsque} \ill{} \ill{} \ill{} \ill{} \ill{} \uncertain{injectifs}, on
\uncertain{définit} une biextension \add{canonique} de $(A, A')$, où $A = G/M$,
$A' = G'/M'$. Dans le cas général, posons
\struck{on dira qu'on a une biextension de}
\[
A = [M \to G], \qquad A' = [M' \to G']
\]
\note{sous $M$ et $G$ : « deg.~$-1$ », « deg.~$0$ » ; sous la ligne, « \ill{} de
chaînes »}

\begin{tikzcd}[column sep=small, row sep=small]
 & A & & & A' & \\
M \arrow[ur, "1"] \arrow[r] & G \arrow[u] & & M' \arrow[ur, "1"] \arrow[r] & G' \arrow[u] &
\end{tikzcd}

\note{deux petits triangles, une flèche marquée $1$ de $M$ vers $A$ ; le sens des
flèches obliques n'est pas sûr}

d'où \ill{} \ill{} \ill{} une \struck{biext} classe canonique
\[
\xi \in \operatorname{Ext}^1(A \overset{L}{\otimes} A', \mathbb{G}_m),
\]
\uncertain{justement} la classe d'une biextension de $(A, A')$ par $\mathbb{G}_m$.
\uncertain{Soit} en effet \ldots\marginal{\uncertain{pour ainsi dire} ?}\note{la note
marginale est écrite en oblique, à gauche d'une accolade qui prend les deux
lignes « justement la classe \ldots{} Soit en effet »}

\struck{Soit}
\[
{}_nA \overset{\text{déf}}{=} \underline{\operatorname{Ext}}^0(\mathbb{Z}/n\mathbb{Z}, A),
\qquad
{}_nA' = \underline{\operatorname{Ext}}^0(\mathbb{Z}/n\mathbb{Z}, A')
\]
\note{l'exposant est lu $0$ ; devant $\underline{\operatorname{Ext}}$, un mot
surchargé}

on \ill{} \ill{}, \ill{} \ill{} le \uncertain{plus haut},

\struck{si}
\[
\begin{cases}
{}_nM = 0 \\
{}_nG = G \quad \text{i.e.}\ nB = B,\ nT = T
\end{cases}
\]
(ce dernier résulte de ${}_nM = 0$)\note{dans l'accolade, après ${}_nM = 0$, un
\struck{${}_nM' = 0$} biffé ; après « \uncertain{i.e.} », un \struck{${}_nB$}
biffé}
\[
0 \to {}_nG \to {}_nA \to M_n \to 0
\]
et $(D(M'_n), {}_nB)$, si $nT = T$ \ldots\ = \struck{\ill{}} \struck{\ill{}}
\[
0 \to {}_nG' \to {}_nA' \to M'_n \to 0
\]
\marginal{donne aussi autre interprétation de ${}_nA$, via hom.\ cobord si ${}_nG =
G$}\note{au-dessus du « $M'$ » biffé, un mot entouré, \ill{} ; « $D(M'_n)$,
${}_nB$ » est surmonté d'une accolade, sous la première suite exacte}

\uncertain{passant} à la limite sur $n = \ell^\nu$, on a
\[
\begin{cases}
0 \to T_\ell(G) \to T_\ell(A) \to M \otimes \mathbb{Z}_\ell \to 0 \\
0 \to T_\ell(G') \to T_\ell(A') \to M' \otimes \mathbb{Z}_\ell \to 0
\end{cases}
\]
\note{le premier terme de la seconde ligne est lu $T_\ell(A')$ sur la page ; on
transcrit $T_\ell(G')$ par symétrie avec la première, sous réserve}

\page{31}
\note{feuillet dactylographié, au verso de la page 30 (reproduit tête-bêche) :
feuillet non paginé du tapuscrit de l'exposé VIII de SGA~7, §2 : (2.2.4) à
(2.2.6), l'accouplement $\ell$-adique associé à $\xi$, la fonctorialité, et 2.3
jusqu'à (2.3.3). Texte publié, non reproduit. Interventions à la main : dans
la marge gauche, entourés, « $\varphi^{(\ell)}$ » et « $\ell^\nu P$ » ; les
$\xi$ en indice ajoutés aux $\phi$ ; « commutatifs » ajouté au-dessus de
« déduits » ; « (pour $n$ fixé) » souligné ; « donnée » au-dessus de « une
compatibilité » ; dans le diagramme (2.3.1), les flèches verticales marquées
« $m$ », « id », « id » ; « (2.3.3) » repassé en marge. Main vraisemblablement
la sienne}

\page{32}

\struck{\ill{} \ill{}} \uncertain{loc.} \ill{} $\mathbb{G}_m$
\[
{}_nA \to A, \qquad {}_nA' \to A'
\]
d'où \uncertain{loc.} \ill{} \uncertain{ces morphismes}

\begin{tikzcd}[column sep=small]
\operatorname{Ext}^1(A \overset{L}{\otimes} A', \mathbb{G}_m) \arrow[r] \arrow[dr] &
\operatorname{Ext}^1({}_nA \overset{L}{\otimes} {}_nA', \mathbb{G}_m) \arrow[d] \\
 & \operatorname{Hom}({}_nA \otimes {}_nA', \mathbb{G}_m)
\end{tikzcd}

\note{le premier $\operatorname{Ext}$ est écrit sur un mot surchargé ; après
$\mathbb{G}_m)$ en haut à droite, un signe \ill{}}

\ill{} \uncertain{vraiment}, dans \ill{} \uncertain{définis} \uncertain{globalt}.
Dans \ill{} \ill{} plus haut \ill{}, donne
\[
\varphi^\xi_n \colon {}_nA \otimes {}_nA' \to \mathbb{G}_m
\]
et par passage à la limite
\[
T_\ell(A) \times T_\ell(A') \longrightarrow T_\ell(\mathbb{G}_m)
\overset{\text{déf}}{=} \mathbb{Z}_\ell(1).
\]

\emph{Prop.}\quad Dans cet accouplement, $D(M'_n)$ est orthogonal à ${}_nG'$
(le \uncertain{reste} si ${}_nM' = 0$) et \struck{$D(M_n)$} ${}_nG$ est orth.\ à
$D(M_n)$ (le \uncertain{reste} si ${}_nM = 0$). Les accouplements induits
\[
{}_nA'/{}_nG' \times D(M'_n) \to \mathbb{G}_m, \qquad
D(M_n) \times {}_nA/{}_nG \to \mathbb{G}_m,
\]
ainsi que ${}_nG/{}_nT \times {}_nG'/{}_nT'$ sont ceux qu'on devine.\note{sous
${}_nA'/{}_nG'$ : « $\simeq M'_n$ » ; sous $D(M'_n)$ : « $= {}_nT'$ » ; sous
$D(M_n)$ : « $= {}_nT$ », sous ${}_nA/{}_nG$ : « $\simeq M_n$ » ; sous ${}_nG/{}_nT$
et ${}_nG'/{}_nT'$ : « $\simeq {}_nB$ », « $\simeq {}_nB'$ ». Une courbe entoure
${}_nA'/{}_nG' \times D(M'_n)$ ; l'énoncé est bordé d'un trait dans la marge}

\emph{Corollaire}.\quad Si l'accouplement entre ${}_nB$ et ${}_nB'$ est une
dualité parfaite (\uncertain{ex.} ${}_nT_n = {}_nT'_n = 0$) \struck{\ill{}} et si
${}_nM = 0$, ${}_nM' = 0$, $nB = B$, $nB' = B'$, alors l'accouplement de ${}_nA$ et
${}_nA'$ est une dualité parfaite.

\page{33}
\note{feuillet dactylographié, au verso de la page 32 (reproduit tête-bêche) :
feuillet non paginé du tapuscrit de l'exposé VIII de SGA~7 — début du §2
« Accouplements définis par une biextension », 2.1 à (2.2.3). Texte publié, non
reproduit. Interventions à la main : « $\varphi^n$ » au-dessus de la flèche
composée de (2.1.2), et les exposants $1$ de $\operatorname{Biext}^1$ ; « (1.3.6) »
corrigé en surcharge ; « $\xi$ » entouré avec « $\in \operatorname{Biext}(P,Q;G)$ »
et relié à « permet » ; « (E) » ajouté au-dessus de « biextension » ; en marge,
entourés, « $\varphi$ » devant (2.1.4) et « $\ell$ » devant (2.2.3) ; les $\xi$
en indice ; « Les relations entre cet accouplement et » biffé et remplacé
par « Il est lié à », « est donné » biffé ; « qu'on laisse au lecteur le soin
d'établir » biffé et remplacé par « cf.~2.3 ci-dessous pour une
démonstration » ; dans le diagramme (2.2.2) les flèches verticales marquées
$m$ ; sous le diagramme, entouré : « où les flèches verticales sont induites
par la multiplication par $m = n'/n$ ». Main vraisemblablement la sienne}

\section*{Picard et dualité}

\note{titre écrit à l'encre et souligné, en haut à droite d'un feuillet par
ailleurs vierge (page 34), qui ne reçoit pas de numéro de page. Pas de
pagination de l'auteur sur les pages 35 à 40}

\page{35}

\note{la page est un tableau à trois colonnes (degrés $0$, $1$, $2$) et trois
lignes (coefficients $\ell$-adiques, $p$-adiques, et $i_p$) ; on le donne ligne
par ligne, colonne par colonne. Le faisceau noté par un $z$ bouclé est rendu
$\mathfrak{z}_\ell$, $\mathfrak{z}_p$, et le groupe noté $U_\ell$ (ou
$\mathfrak{U}_\ell$) est rendu $U_\ell$. Une accolade embrasse les deux
premières lignes, avec, dans la marge, écrit de bas en haut : « et
\uncertain{cohomologie discrète} »}

\emph{Première ligne} ($\ell$) :
\[
H^0(X, \mathfrak{z}_\ell) = U_\ell,
\]
\[
\begin{aligned}
H^2(X, \mathbb{Z}/\ell) &= H^2(X, \mathfrak{z}_\ell) \otimes_{\mathbb{Z}/\ell}
\operatorname{Hom}(U_\ell, \mathbb{Z}/\ell) \\
&= \mathbb{Z}/\ell \otimes_{\mathbb{Z}/\ell}
\operatorname{Hom}(U_\ell, \mathbb{Z}/\ell) \simeq \operatorname{Hom}(U_\ell, \mathbb{Z}/\ell) ;
\end{aligned}
\]
\note{dans le dernier membre, un $\mathbb{Z}/\ell$ est biffé devant
$\operatorname{Hom}$}
$H^1(X, \mathfrak{z}_\ell) = P^{(\ell)}$ \uncertain{noyau de la} jacobienne $P$ de $X$
\uncertain{annulé par} $\ell$ ;
\[
H^1(X, \mathbb{Z}/\ell) = H^1(X, \mathfrak{z}_\ell) \otimes_{\mathbb{Z}/\ell}
\operatorname{Hom}(U_\ell, \mathbb{Z}/\ell)
= P^{(\ell)} \otimes_{\mathbb{Z}/\ell} \operatorname{Hom}(U_\ell, \mathbb{Z}/\ell) ;
\]
\note{entre les deux égalités, une première écriture « $= P^{(\ell)} \otimes
\operatorname{Hom}(\mathfrak{z}_\ell, H)$ » est biffée}
or $P^{(\ell)} \times P^{(\ell)} \to U_\ell$ \ill{} \ill{}
\marginal{\ill{} l'autodualité de $P$ \ldots{} et la dualité de \uncertain{Kummer
Serre} \ldots}
\note{la note entre crochets est encadrée, sous l'accouplement}
;
\[
H^2(X, \mathfrak{z}_\ell) \simeq \mathbb{Z}/\ell, \qquad
H^0(X, \mathbb{Z}/\ell) \simeq \mathbb{Z}/\ell .
\]

\emph{Deuxième ligne} ($p$) :
\[
H^0(X, \mathfrak{z}_p) = 0, \qquad
H^2(X, \mathbb{Z}/p) \simeq H^2(X, \mathbb{Z}/p) = 0 ;
\]
$H^1(X, \mathfrak{z}_p) = P^{(p)}$ : \uncertain{noyau de $p$} de la jacobienne ;
$H^1(X, \mathbb{Z}/p) \simeq$ \uncertain{noyau} de $H^1(X, \underline{\mathcal{O}}_X)$
\uncertain{fixes} par $F$
\marginal{dualité de Serre-Cartier}
\note{encadré sous la case}
;
\[
H^2(X, \mathfrak{z}_p) \simeq \mathbb{Z}/p, \qquad H^0(X, \mathbb{Z}/p) \simeq
\mathbb{Z}/p .
\]
\note{« $H^2(X, \mathbb{Z}/p) \simeq H^2(X, \mathbb{Z}/p)$ » est tel quel sur la
page}

\emph{Troisième ligne} ($i_p$) :
\[
H^0(X, i_p) = 0, \qquad H^1(X, i_p) \simeq {}_F H^1(X, \underline{\mathcal{O}}_X),
\qquad H^2(X, i_p) = H^1(X, \underline{\mathcal{O}}_X)_F .
\]
\note{le coefficient de la case du milieu est surchargé ($\mathfrak{z}_p$ corrigé
en $i_p$ ?) ; les indices $F$ sont à gauche en bas pour le noyau, à droite en
bas pour le conoyau, tels qu'écrits}

\page{36}

\[
H^1(X, \underline{\operatorname{Ext}}^1(f^{*}(B), \mathbb{G}_{m\,X}))
\]
\note{une accolade sous $\underline{\operatorname{Ext}}^1(f^{*}(B), \mathbb{G}_{m\,X})$
renvoie à la ligne suivante}
\[
\operatorname{Ext}_X(f^{*}(B), \mathbb{G}_{m\,X}) \simeq
\operatorname{Ext}_S(B, f^{!}(\mathbb{G}_{m\,X}))
\]
\note{au-dessus de $f^{*}(B)$ est écrit $A$ ; la ligne suivante est reliée à
celle-ci par une double flèche verticale $\Uparrow$}
\[
H^p(X, \underline{\operatorname{Ext}}^q_X(A, \mathbb{G}_{m\,X}))
\]
\note{un trait horizontal sépare ce calcul de la liste qui suit}

\begin{itemize}
\item \uncertain{Suites exactes de cohomologie}. Critères de
\uncertain{representabilité} ;
\item \uncertain{Frobenius} \uncertain{Verschiebung} ;
\item \uncertain{Dualité de Cartier} [\uncertain{\ill{} Pic \ldots}] ;
\item \uncertain{Th.\ de} \uncertain{finitude} du V.A.\ (\uncertain{Th.\ de
\ill{}}) ;
\item \uncertain{Propriétés formelles des groupes} \ldots\ ;
\item \uncertain{Propriétés différentielles des $\ldots$ groupes et fibrés} ;
\item \uncertain{Cas particulier} : \emph{fibrés linéaires}, symplectiques,
orthogonaux etc.\ \ldots
\end{itemize}
\note{liste en accolade, écrite rapidement ; presque chaque ligne est
incertaine. La sixième ligne a un premier mot surchargé ; « fibrés linéaires »
est souligné deux fois}

\uncertain{Sous \ill{} pour la \ill{}} \ill{} ??

\[
\underline{\operatorname{Hom}}(\mathbb{G}_a, \mathbb{G}_m) \qquad
\mathbb{Z} \to \mathbb{Z}
\]
\note{au pied de la page, sous un second trait : $\underline{\operatorname{Hom}}(\mathbb{G}_a,
\mathbb{G}_m)$ semble biffé d'un trait ; « $\mathbb{Z} \to \mathbb{Z}$ » est lu
sous réserve}

\page{38}

$X$ propre et simple sur $k$ [\uncertain{généralisations} à une base
quelconque]

$A$ V.A.\ sur $k$
\[
H^1(X, A) = \varinjlim_n {}_nH^1(X, A) ;
\]
d'autre part \ill{} \ill{} une suite exacte
\[
0 \to \operatorname{Hom}(X, A)_n \to H^1(X, {}_nA) \to {}_nH^1(X, A) \to 0
\]
\uncertain{qui s'écrit} aussi
\[
0 \to \operatorname{Hom}(\operatorname{Alb}_X, A)_n \to
\operatorname{Hom}(D({}_nA), \operatorname{Pic}_X) \to {}_nH^1(X, A) \to 0
\]
\note{devant $H^1(X, {}_nA)$, un signe biffé}
\uncertain{or} \uncertain{utilisant} \ill{}
\[
\operatorname{Hom}(\operatorname{Alb}_X, A) \simeq \operatorname{Hom}(A',
(\operatorname{Alb}_X)') = \operatorname{Hom}(A', \operatorname{Pic}^0_X),
\qquad D({}_nA) \simeq {}_nA'
\]
\uncertain{une} suite exacte
\[
(*) \qquad 0 \to \operatorname{Hom}(A', \operatorname{Pic}^0_X)_n \to
\operatorname{Hom}({}_nA', \operatorname{Pic}_X) \to {}_nH^1(X, A) \to 0
\]
et passant à la limite \add{\uncertain{inductive}} sur $n$ :
\[
\boxed{0 \to \operatorname{Hom}(A', \underline{\operatorname{Pic}}^0_X)
\otimes_{\mathbb{Z}} \mathbb{Q} \to \operatorname{Hom}(T_{\cdot}(A'),
\underline{\operatorname{Pic}}_X) \to H^1(X, A) \to 0}
\]
\note{la formule est encadrée ; « inductive » est écrit au-dessus de « sur », et
l'exposant $0$ de $\operatorname{Pic}$ est surchargé. Le « $T_{\cdot}$ » est
écrit ainsi, avec un point en indice}

\uncertain{Or} la suite exacte $(*)$ \uncertain{se compare} à la suite exacte
\uncertain{déduite} de $0 \to {}_nA' \to A' \to A' \to 0$ \uncertain{en
appliquant} \struck{$\operatorname{Ext}$} $\underline{\operatorname{Hom}}(-,
\operatorname{Pic}_X)$, \ldots\ \uncertain{laquelle}
\[
(**) \qquad 0 \to \operatorname{Hom}(A', \underline{\operatorname{Pic}}_X)_n \to
\operatorname{Hom}({}_nA', \operatorname{Pic}_X) \to {}_n\operatorname{Ext}^1(A',
\underline{\operatorname{Pic}}^0_X) \to 0
\]
\uncertain{qui donne} à la limite
\[
\boxed{0 \to \operatorname{Hom}(A', \operatorname{Pic}_X) \otimes_{\mathbb{Z}}
\mathbb{Q} \to \operatorname{Hom}(T_{\cdot}(A'), \operatorname{Pic}_X) \to
\operatorname{Ext}^1(A', \underline{\operatorname{Pic}}_X) \to 0}
\]
\note{l'exposant de $\operatorname{Pic}$ dans le dernier terme de $(**)$ est lu
$0$ sous réserve ; la formule finale est encadrée et bute sur le bord droit du
feuillet}

\page{39}

\uncertain{Donc} \uncertain{on en conclut un} \uncertain{isom.} canonique
\[
\boxed{H^1(X, A) \simeq \operatorname{Ext}^1(A', \underline{\operatorname{Pic}}_X)}
\]
($A$, variété abélienne), \uncertain{qu'on comparera} à \uncertain{la}
\uncertain{formule} \uncertain{de dualité}
\[
\boxed{H^1(X, G) \simeq \operatorname{Hom}(D(G), \underline{\operatorname{Pic}}_X)}
\]
pour $G$ commutatif fini.
\marginal{Cet isomorphisme \uncertain{reste-t-il} valable pour $X$ singulier ???
Il est possible qu'on ait besoin de $H^i(X, \mathbb{G}_m) = 0$ pour $i \geqslant 2$,
qui \uncertain{est faux} \ill{} $X$ \uncertain{non sing.}??}
\note{la note marginale est écrite à gauche du premier encadré, sur neuf lignes
courtes ; le « $\simeq$ » du premier encadré est une double demi-flèche}

\uncertain{Remarque}, \uncertain{on a des} homs canoniques

\begin{tikzcd}[row sep=small]
\operatorname{Ext}^1(\underline{\operatorname{Alb}}_X, A) \arrow[r] \arrow[d, "\wr"] & H^1(X, A) \\
\operatorname{Ext}^1(A', \underline{\operatorname{Alb}}_X{}') \arrow[d, "\wr"] & \\
\operatorname{Ext}^1(A', \underline{\operatorname{Pic}}^{\tau}) &
\end{tikzcd}

\uncertain{dont on vérifie} \uncertain{que ce sont} des \emph{isomorphismes}
\emph{dans} \ill{}, qui définissent sur les \ill{} \uncertain{groupes} \ill{}
\uncertain{une partie} \uncertain{l'isomorphisme} \uncertain{envisagé plus haut}.
\note{l'exposant de $\underline{\operatorname{Pic}}$ dans la dernière ligne du
diagramme est lu $\tau$, sous réserve ; « isomorphismes » et « dans » sont
soulignés ; les trois dernières lignes sont une écriture rapide, lue en partie
seulement}

\page{40}

\note{à gauche du premier encadré, une flèche verticale $X \to S$. On rend le
faisceau de cohomologie relative, qu'il note d'un $\mathcal{H}$ cursif souligné,
par $\underline{\mathcal{H}}$, et le groupe noté par un $\gamma$ bouclé par
$\gamma$}
\[
\boxed{\underline{\mathcal{H}}^1(X/S; \underline{\operatorname{Hom}}_{S\text{-gr}}(\gamma,
G)) \simeq \underline{\operatorname{Hom}}_{S\text{-gr}}(\gamma,
\underline{\mathcal{H}}^1(X/S; G))}
\]
isom de foncteurs, \struck{\ill{}} quand tout est \emph{affine} sur $S$ (?)
\[
\boxed{\underline{\mathcal{H}}^1(X/S; \underline{\operatorname{Ext}}^1_{S\text{-gr}}(\gamma,
G)) \simeq \underline{\operatorname{Ext}}^1_{S\text{-gr}}(\gamma,
\underline{\mathcal{H}}^1(X/S; G))}
\]
Sous quelles conditions ???

\note{puis, sous $P_X$, un petit tableau à deux lignes : sur $X$,
« $\gamma_X$ » et « $\gamma_X \times \gamma_X$ » (tous deux surchargés), sur $S$,
« $\gamma$ » et « $\gamma \times \gamma$ », $X$ et $S$ reliés par un trait
vertical}

\[
\underline{\mathcal{H}}^1(X/S, \underline{\operatorname{Hom}}_{S\text{-gr}}(\gamma, G))
\to \underline{\operatorname{Hom}}(\gamma, \underline{\mathcal{H}}^1(X/S, G))
\]
\[
\underline{\mathcal{H}}^1(X, \operatorname{Hom}(A, B)) \simeq
\]
\struck{\ill{}}
\[
\boxed{\begin{array}{c}
H^1(X/S; \underline{\operatorname{Hom}}_{S\text{-gr}}(\gamma, G)) \\
\downarrow \\
\underline{\operatorname{Hom}}_{S\text{-gr}}(\gamma, H^1(X/S, G))
\end{array}}
\]
\note{le dernier encadré est dessiné avec une colonne vide à droite ; la
ligne « $\underline{\mathcal{H}}^1(X, \operatorname{Hom}(A, B)) \simeq$ » reste
sans second membre}

\end{document}
